% Tapis de Quillen, with Quillen's offprint and preprint (cote 162-5, typed part): archive cote 162-5, site page 45 (PDF page 46).
% typescript. Draft compared: none.
% AI-transcribed by Claude Opus 5.5 (claude-opus-5-5, adaptive thinking, effort high), Message Batches API, batch t162_full (msgbatch_01UAh3zyWPd454AeDkhnmQ33), 2026-11-13; not yet checked by eye.
\scanpage{162-5/45}
g manifestement se prolonge en $\overline{g} : \overline{T} \to X\mathrm{x}Y$ . Cela nous ramène à prouver
\inserted[typed]{le composé}
que \hand{\unsure{$\Gamma$}}$\overline{T}'$\hand{$\xrightarrow{\alpha_i}$}\hand{\unsure{$\Gamma$}}$Z_i^!$\hand{$\xrightarrow{\alpha_{i*}}$}$\overline{T}'$ ne dépend pas de i=0,1. Ecrivons X au lieu de $\overline{T}'$,et
considérons le diagramme cartésien d'inclusions
\diagram{\begin{tikzcd}
X_s \arrow[r, "\hand{\alpha_s}"] \arrow[d, "\hand{\alpha_s}"'] & X \arrow[d, "\hand{\Gamma = (p,\mathrm{id}_X)}"] \\
X \arrow[r, "i_s"] & S^1\mathrm{x}X
\end{tikzcd}}
,

on a alors $\hand{\alpha}^{\hand{*}}_s \hand{\alpha}_{s*} = \Gamma^* i_{s\hand{*}}$ , qui ne dépend pas de s puisque $i_s$ n'en
dépend pas à homotopie près. - Le reste est sans doute l'AQT.

Sur la strcture des B.(X). Ce sont tous des algèbres graduées
sur B.(pt)=\hand{$\Lambda$} , qui est une algèbre graduée à degrés positifs, réduite
\margin{\unsure{\S}}
à $\underline{F}_2$ en degré\xout{s} 0. C'est d'ailleurs une algèbre de polynômes \inserted{(à une infinité de générateurs \unsure{$x_i$ ($x_i$ de degré de la forme $2^i-1$)}} \deleted{\unsure{connue}} ...
\inserted[typed]{d'\unsure{anneaux gradués}}
On a un homomorphisme évident B.(X) $\to$ H$^{\over{\cdot}{\unsure{*}}}$(X), nul en degrés \hand{$>$}0, qui
à tout Z sur X associe l'image de 1\hand{$\in$}H$^0$(Z) par l'homomorphisme de Gysin
associé à Z$\to$ X. Cet homomorphisme par Thom est surjectif (toute clas-%
\margin{en fait $H(X) \simeq B.(X)/\Lambda^+B.(X)$}
se de cohomologie provient d'une image de variété compacte ...); \xout{noter}
On ppouve même qu'il existe un relèvement fonctoriel en X
\margin{?}
H$^\bullet$(X) $\to$ B.(X), \inserted{il est douteux qu'il soit}\hand{(}compatible avec multiplication\hand{)} \xout{tel}tel que l'on ait\hand{:}

\over{a)}{\unsure{1)}} L'homomorphisme correspondant H$^\bullet$(X)$\otimes_{\underline{F}_2}$\hand{$\Lambda$} \hand{$\simeq$}B.(X) est un iso-%
\margin{?}
morphisme (ce sera vrai pour tout relèvement individuellement ...), \deleted{et}

\over{b)}{\unsure{2)}} \hand{[}\inserted{Quillen ignore si} \over{Le\xout{x}}{\unsure{le}} relèvement est compatible aussi avec les hom. de Gysin.\hand{]}

Ainsi, la catégorie \underline{B} admet un foncteur pleinement fidèle dans la
catégorie des modules gradués libres de type fini sur \hand{$\Lambda$}, les modules
de type fini qu'on obtient ainsi ayant des générateurs à degrés \hand{$<$} 0 et
satisfaisant (pour X connexe i.e. quand il y a un seul générateur de
degré zéro\hand{)} la dualité de Poincaré. Quand on prend la catég\unsure{o|b}rie quasi-%
abélienne\hand{-}\inserted{\unsure{B}}associée à \underline{B}, on trouve donc tous les modules gradués libres
de type fini à générateurs \xout{négatifs} en degrés \hand{$\leq$} 0 seulement. C'est une

\mspage{45}
